%% TEMPLATE — ltx-talk shared preamble (adapt the PLACEHOLDERS marked <<...>>).
%% Copy this next to your decks (e.g. as ltx-common.tex) and fill in the
%% Identity and ThemeAccent blocks below; everything else adapts automatically.
%% Load it AFTER \documentclass{ltx-talk}. \DocumentMetadata must already be set
%% (it must be the very first thing in the file, before \documentclass) so
%% tagging is active.
%%
%% Design notes:
%%  - ltx-talk has no \usetheme / \setbeamercolor / \setbeamertemplate. The
%%    visual style is driven by the kernel template system via \EditInstance.

%% ---------------------------------------------------------------------------
%% Provenance: which ltx-talk and which converter this course was built against
%% is recorded in ltx-talk-conversion.toml at the root of this repo, written by
%% scripts/stamp_conversion.py. Deliberately not repeated here -- a value copied
%% into this file would have to be maintained by hand, and this template's own
%% <<placeholder>> markers are the standing evidence for how that goes.
%% ---------------------------------------------------------------------------

%% ---------------------------------------------------------------------------
%% Accessibility: one line that a PDF/UA checker will otherwise reject every
%% deck in the course over. Pure tagging — no visual change at all.
%% See references/compromises.md A-HEADINGS.
%%
%%  * ltx-talk hard-codes frametitle -> H4, while \section is H1, so the first
%%    heading in a deck is an orphan H4 ("headings do not begin at level one").
%%    H2 sits directly under the section H1.
%% ---------------------------------------------------------------------------
\tagpdfsetup{role/new-tag = frametitle / H2}

%% Do NOT add \tagpdfsetup{math/alt/use} here. It puts an /Alt on every Formula
%% element, and a screen reader that finds /Alt reads it INSTEAD of the MathML,
%% which is the accessible representation of maths under PDF/UA-2. The kernel
%% leaves this switch off for ua-2 deliberately (it is auto-on for ua-1 only) --
%% that is the correct behaviour, not an oversight to work around. A validator
%% that demands /Alt on Formula is testing against ua-1. See A-MATHALT.

%% ---------------------------------------------------------------------------
%% Colours
%%
%% Emphasis colours are darkened for WCAG AA (>= 4.5:1 on white); the saturated
%% originals look fine on a projector and fail a checker. Ratios in comments.
%% If you add your own, measure it — see compromises.md A-CONTRAST for the
%% formula and the >=200 dpi caveat.
%% ---------------------------------------------------------------------------
\usepackage{xcolor}
\definecolor{cream}{rgb}{1,0.99,0.89}
\definecolor{myellow}{rgb}{0.95,0.33,0}
% \definecolor{green}{rgb}{0.2,0.8,0}                    % 2.15:1 — fails AA
\definecolor{green}{rgb}{0,0.50,0}                       % 5.17:1
\definecolor{blue}{rgb}{0.0039,0.274,1}                  % 6.33:1
%% Decks write {\color{red}...} inline far more often than they use a macro.
%% Redefining the name once fixes every site — but ONLY if nothing uses red
%% graphically. Check first:
%%   grep -hoE '(draw|fill|text|color)\s*=\s*red[^,;}]*|red![0-9]+' *.tex
\definecolor{red}{rgb}{0.75,0,0}                         % was 1,0,0 = 4.00:1
\definecolor{aired}{rgb}{0.75,0,0}                       % 6.52:1
\definecolor{aiteal}{rgb}{0.05,0.40,0.40}                % 6.75:1 (was teal!80)
\definecolor{aiorange}{rgb}{0.68,0.33,0}                 % 5.17:1 (was orange)
\definecolor{aigreen}{rgb}{0,0.50,0}                     % 5.17:1 (was dkgreen)
\definecolor{lightblue}{rgb}{0.8,1,1}
\definecolor{yellow}{rgb}{1,0.99,0.3}
\definecolor{BoxCol}{rgb}{0.95,0.33,0}
\definecolor{brightmaroon}{rgb}{0.76, 0.13, 0.28}
\definecolor{bostonuniversityred}{rgb}{0.8, 0.0, 0.0}
\definecolor{aliceblue}{rgb}{0.94, 0.97, 1.0}
\definecolor{cadmiumgreen}{rgb}{0.0, 0.42, 0.24}
\definecolor{darkgreen}{rgb}{0.01, 0.75, 0.24}
\definecolor{darkorange}{rgb}{0.91, 0.41, 0.17}
\definecolor{violet}{rgb}{0.56, 0.0, 1.0}

%% ---------------------------------------------------------------------------
%% Theme accent colours — EDIT THESE TWO LINES to match your beamer theme.
%%
%% Common beamer theme → ltx-talk colour mapping:
%%   CambridgeUS / Madrid  →  \ThemeAccent = bostonuniversityred   \ThemeAccentText = white
%%   Warsaw / Frankfurt    →  \ThemeAccent = {rgb}{0.2,0.2,0.6}    \ThemeAccentText = white
%%   Metropolis            →  \ThemeAccent = {rgb}{0.35,0.35,0.35}  \ThemeAccentText = white
%%   AnnArbor              →  \ThemeAccent = {rgb}{0.0,0.32,0.59}   \ThemeAccentText = {rgb}{0.96,0.76,0.0}
%%   Boadilla / Rochester  →  \ThemeAccent = {rgb}{0.2,0.4,0.65}    \ThemeAccentText = white
%%   Luebeck / Copenhagen  →  \ThemeAccent = {rgb}{0.1,0.1,0.5}     \ThemeAccentText = white
%%   Singapore             →  \ThemeAccent = {rgb}{0.0,0.5,0.5}     \ThemeAccentText = white
%%
%% You can also use any named colour defined above (e.g. brightmaroon,
%% cadmiumgreen) or add a new \definecolor before these lines.
%% ---------------------------------------------------------------------------
\newcommand{\ThemeAccent}{bostonuniversityred}   %% <<ADAPT>> header/footer bg
\newcommand{\ThemeAccentText}{white}             %% <<ADAPT>> header/footer fg

%% ---------------------------------------------------------------------------
%% Identity — EDIT THESE FIVE LINES, everything else auto-adapts.
%% ---------------------------------------------------------------------------
\newcommand{\AuthorShort}{<<SHORT AUTHOR(S)>>}        %% e.g. Smith and Jones
\newcommand{\AuthorLong} {<<FULL AUTHOR NAME(S)>>}    %% e.g. Alice Smith \and Bob Jones
\newcommand{\AuthorEmails}{<<EMAIL(S)>>}              %% e.g. a.smith@uni.ac.uk \quad b.jones@uni.ac.uk
\newcommand{\InstShort}  {<<SHORT INSTITUTION>>}      %% e.g. Aberdeen
\newcommand{\InstLong}   {<<FULL INSTITUTION NAME>>}  %% e.g. University of Aberdeen

%% Shaped exactly as the beamer source writes it, so the deck round-trips: the
%% emails sit INSIDE \author, after a \\, and \and separates co-authors. Both
%% work verbatim on the ltx-talk title page -- the titlepage template rebinds
%% \and to \quad itself (C-AND-TITLE), and \\ is an ordinary line break there.
\author[\AuthorShort]{\AuthorLong\\ \AuthorEmails}
\institute[\InstShort]{\InstLong}

%% ---------------------------------------------------------------------------
%% Content styling (uses \ThemeAccent — change the macro above, not here)
%% ---------------------------------------------------------------------------

%% Coloured, bold \emph. Individual decks may override with \let\emph\textbf.
\renewcommand{\emph}[1]{{\color{\ThemeAccent}\textbf{#1}}}

\newcommand{\obl}[0]{\mbox{\color{violet}{$\mathbf{O}$}}}
\newcommand{\prh}[0]{\mbox{\color{red}{$\mathbf{F}$}}}
\newcommand{\per}[0]{\mbox{\color{green}{$\mathbf{P}$}}}
\newcommand{\viz}[0]{{\it viz.},\ }

%% ---------------------------------------------------------------------------
%% Maths and content packages
%% ---------------------------------------------------------------------------
\usepackage{bm}
\renewcommand{\mathbf}[1]{\ensuremath{\bm{#1}}}

\usepackage{color,ulem}
\usepackage{algorithm}
\usepackage{bbm}
\usepackage{amsmath,amssymb}
%% Pseudocode: use the classic algorithmicx/algpseudocode engine, NOT
%% algpseudocodex. Under ltx-talk's tagged output, algpseudocodex's algorithmic
%% body triggers "Improper \halign inside $$'s" whenever \Comment/\Function are
%% used (it cannot be typeset at all). The classic engine is tagging-safe.
%% [noend] reproduces algpseudocodex's noEnd look (hides \EndFunction etc.).
\usepackage{algorithmicx}
\usepackage[noend]{algpseudocode}
\usepackage{wasysym}
\usepackage{multirow}
\usepackage{mathtools}
\usepackage{xifthen}
\usepackage{cancel}
\usepackage{pifont}
\newcommand{\cmark}{\ding{51}}%
\newcommand{\xmark}{\ding{55}}%

\usepackage{graphicx}

\usepackage{pgfplots}
\pgfplotsset{width=7cm,compat=1.15}

%% Speaker-note stubs (no-op), kept for source compatibility with pdfpc/beamer.
\newcommand{\pdfpcnote}[1]{}
\newcommand{\note}[2][]{}    %% beamer \note[item]{text} / \note{text}

%% ---------------------------------------------------------------------------
%% alertblock / exampleblock — beamer's coloured highlight boxes.
%%
%% UNLIKE columns/column/block below, these are NOT native in ltx-talk 0.5.3
%% (only `block' is — see C-NATIVE-ENVS). ltx-talk has no built-in
%% alertblock/exampleblock at all, so this stub stays ACTIVE. See C-ALERTBLOCK.
%%
%% \begin{alertblock}{Title}...\end{alertblock} — mandatory brace argument,
%% same call syntax as beamer's own. Do NOT try to preserve beamer's optional
%% overlay spec (`\begin{alertblock}<2->{Title}`) with a d<>m/s<>m wrapper that
%% forwards the title as `[{#2}]` — tcolorbox parses `[...]` as its key=value
%% OPTIONS argument, not a positional title, so the box title silently becomes
%% the literal character `[`. If a deck needs the overlay behaviour, wrap the
%% whole environment in `\onslide<2->{\begin{alertblock}{Title}...}` instead of
%% inventing a `<>`-swallowing wrapper.
%%
%% Accessibility note: the box title (first argument) is tagged as plain text,
%% not as a heading; screen readers see it as body text. If semantic structure
%% matters, use a \subsubsection* inside the box body instead. The coloured
%% border is drawn by tcolorbox's standard (non-TikZ) skin and is automatically
%% marked as an Artifact by the LaTeX tagging kernel — it does not pollute the
%% reading order.
%%
%% We load tcolorbox WITHOUT [most]: the skins/TikZ libraries in [most] produce
%% untagged graphics that can break the PDF tag tree under \DocumentMetadata.
\usepackage{tcolorbox}

\newtcolorbox{alertblock}[1]{
  title={#1}, fonttitle=\bfseries,
  colback=red!8!white, colframe=red!70!black,
  left=4pt, right=4pt, top=2pt, bottom=2pt,
}
\newtcolorbox{exampleblock}[1]{
  title={#1}, fonttitle=\bfseries,
  colback=cadmiumgreen!8!white, colframe=cadmiumgreen,
  left=4pt, right=4pt, top=2pt, bottom=2pt,
}

%% ###########################################################################
%% ##  STOP.  DO NOT USE THE STUBS BELOW WITH ltx-talk.                     ##
%% ##                                                                        ##
%% ##  ltx-talk provides columns / column / block / frame* NATIVELY          ##
%% ##  (ltx-talk.cls lines 1074, 1141, 2134, 997). Defining these stubs on   ##
%% ##  top of it CLASHES ("Command \columns already defined").               ##
%% ##  See compromises.md -> C-NATIVE-ENVS.                                  ##
%% ##                                                                        ##
%% ##  Use the native environments. The stubs are kept only for a plain      ##
%% ##  kernel-class setting where they genuinely do not exist.               ##
%% ##                                                                        ##
%% ##  If you DO need theorem-like environments, ltx-talk has none: build    ##
%% ##  them with tcolorbox (C-THEOREM), NOT with the native \block, which    ##
%% ##  collides with algorithmicx's csname stack (C-BLOCK-ALGO).             ##
%% ###########################################################################
\iffalse   %% <-- stubs disabled; delete this \iffalse...\fi only for a non-ltx-talk class

%% ---------------------------------------------------------------------------
%% Beamer-compatibility stubs — constructs beamer provided that ltx-talk does
%% not offer natively. See also: compromises.md entries COLS-01 and BOX-01.
%% ---------------------------------------------------------------------------

%% columns / column — beamer's side-by-side layout.
%% Replaced by minipage pairs. Accessibility note: the PDF tag tree records
%% content in source order, so screen readers hear the ENTIRE left column
%% before the right column. This is correct for independent side-by-side
%% blocks (Problem | Solution, Code | Output) — the normal use of beamer
%% columns. Do NOT use this for content where the two columns have paired rows
%% (e.g. a concept grid); use a proper tabular/table instead, whose rows are
%% read left-to-right as expected. The \hfill between columns is decorative
%% and does not appear in the tag tree. The optional alignment arg on
%% `columns' is silently accepted but ignored (minipage always uses [t]).
\newenvironment{columns}[1][t]{\par\medskip\noindent\ignorespaces}{\par\medskip}
\newenvironment{column}[2][t]{\begin{minipage}[#1]{#2}\ignorespaces}%
                              {\end{minipage}\hfill\ignorespaces}

%% block — beamer's plain highlight box. NATIVE in ltx-talk; this stub is only
%% for a plain kernel-class setting. See the alertblock/exampleblock stub
%% ABOVE this disabled section for the two boxes that DO still need defining.
\newtcolorbox{block}[1]{
  title={#1}, fonttitle=\bfseries,
  colback=aliceblue, colframe=\ThemeAccent,
  left=4pt, right=4pt, top=2pt, bottom=2pt,
}

\fi   %% <-- end of the disabled Beamer-compatibility stubs (see C-NATIVE-ENVS)

%% ---------------------------------------------------------------------------
%% Things every real ltx-talk conversion turns out to need. Unlike the stubs
%% above, these are ACTIVE — each is a compromises.md entry paid for in blood.
%% ---------------------------------------------------------------------------

%% C-FRAMESTAR-TAG: ltx-talk's frame* re-tokenises its body, which under
%% \DocumentMetadata corrupts the tag tree as soon as a listing appears in it
%% (70 errors on one real deck). Suspend tagging around the WHOLE frame*.
%% Use \tag_stop:/\tag_start:, NOT \SuspendTagging (which leaves a dangling
%% marked-content unit and breaks again on adjacent frame*).
%% Trade-off: verbatim slides become artifacts, i.e. not in the reading order.
\ExplSyntaxOn
\AddToHook{env/frame*/before}{\tag_stop:}
\AddToHook{env/frame*/after}{\tag_start:}
%% C-ALGO-FLOAT: the `algorithm' float is not registered with the tagging float
%% module. (Better still: drop the float and keep bare, tagged `algorithmic'.)
\AddToHook{env/algorithm/before}{\tag_stop:}
\AddToHook{env/algorithm/after}{\tag_start:}
\ExplSyntaxOff

%% C-CALL-NEST: classic \Call is \textproc{#1}\ifthenelse{\equal{#2}{}}{}{(#2)},
%% and \equal cannot cope with a nested \Call in #2. Always emit parentheses.
\algrenewcommand\Call[2]{\textproc{#1}(#2)}

%% A-TIKZ-ALT: an \input of a generated figure (xfig .pdf_t, .pspdftex, .pgf)
%% produces NO warning and lands with no usable /Alt -- it is missing from the
%% alt-text debt rather than listed as missing, because the \includegraphics is
%% in a separate, generated file that no linter over the deck sources ever sees.
%%
%% For a tikzpicture or a pgfplots axis you do NOT need any of this: use the
%% alt key on the environment itself, which latex-lab provides --
%%   \begin{tikzpicture}[alt=A square joined to a circle by an arrow.]
%%
%% For the inputted-figure case there is no such option, so set the same
%% latex-lab key at the call site and let the \includegraphics inside pick it up:
%%   \altinput{A generated diagram showing a labelled box.}{fig.pdf_t}
%% Do not edit the .pdf_t to add alt= -- xfig regenerates it.
\ExplSyntaxOn
\NewDocumentCommand{\altinput}{ m m }
  { \group_begin: \keys_set:nn { tag / graphic } { alt = {#1} } \input{#2} \group_end: }
\ExplSyntaxOff

%% C-FRAMESUBTITLE: \framesubtitle is accepted by ltx-talk 0.5.3 and typesets
%% NOTHING -- the class declares, sets and clears \g__talk_frame_subtitle_tl but
%% never reads it. No error, no warning, and the page count is unchanged, so the
%% text simply vanishes and every fidelity check still passes. This is a declared
%% upstream gap, not a bug: ltx-talk.pdf says the frame subtitle "is not used in
%% output: this will be addressed in later releases."
%% Fold both parts into the one thing the class does typeset, the frame title:
%%
%%   \frametitlesub{<title>}{<subtitle>}
%%
%% Defined as a NEW macro rather than a redefinition of \framesubtitle on
%% purpose: appending to the title needs either the private
%% \g__talk_frame_title_tl or a wrapper mirroring \frametitle's `D<>{all} O{#3} m'
%% signature, and the same manual warns that "the final form of the frame title
%% interface is not decided". Do not couple decks to that.
%% Keep this in sync anywhere a deck might see it (cf. altfigure above): a deck
%% that still dual-compiles under plain Beamer wants the SAME macro name defined
%% there as a real \frametitle + \framesubtitle pair --
%%   \newcommand{\frametitlesub}[2]{\frametitle{#1}\framesubtitle{#2}}
%% -- so the source stays exportable both ways.
%% Retire this once a later ltx-talk release typesets the subtitle natively.
\newcommand{\frametitlesub}[2]{\frametitle{#1 --- {\normalsize #2}}}

%% C-OLDFONT: ltx-talk does not define the obsolete two-letter font commands,
%% and decks use them inside maths too, where a text-shape switch is illegal.
\providecommand{\sc}{\ifmmode\else\scshape\fi}
\providecommand{\it}{\ifmmode\else\itshape\fi}
\providecommand{\bf}{\ifmmode\else\bfseries\fi}
\providecommand{\rm}{\ifmmode\else\rmfamily\fi}
\providecommand{\sf}{\ifmmode\else\sffamily\fi}
\providecommand{\tt}{\ifmmode\else\ttfamily\fi}
\providecommand{\sl}{\ifmmode\else\slshape\fi}

%% C-FONTS: ltx-talk's pdfLaTeX path loads sansmathfonts + lmodern[nomath] and
%% sets \rmdefault=\sfdefault, so MATHS is sans. Uncomment to restore serif
%% maths with sans body text (i.e. beamer's \usefonttheme[onlymath]{serif}).
%% Must come after amssymb.
% \DeclareSymbolFont{operators}   {OT1}{lmr} {m}{n}
% \DeclareSymbolFont{letters}     {OML}{lmm} {m}{it}
% \DeclareSymbolFont{symbols}     {OMS}{lmsy}{m}{n}
% \DeclareSymbolFont{largesymbols}{OMX}{lmex}{m}{n}
% \SetSymbolFont{operators}{bold}{OT1}{lmr} {bx}{n}
% \SetSymbolFont{letters}  {bold}{OML}{lmm} {b}{it}
% \SetSymbolFont{symbols}  {bold}{OMS}{lmsy}{b}{n}

%% hyperref metadata. With \DocumentMetadata active hyperref is already loaded;
%% we only adjust its settings here.
%%
%% pdfauthor is NOT decoration. ltx-talk copies \author's argument into the
%% document metadata (\g__tag_title_author_tl -> XMP dc:creator), so the emails
%% and \\ that the beamer source keeps inside \author land in dc:creator too --
%% measured on 0.6.2: dc:creator reads "A. Lecturer  a.l@uni.ac.uk". Setting
%% pdfauthor here wins, and fixes it for every deck without touching a single
%% deck's \author line. See references/compromises.md C-TITLEPAGE.
\hypersetup{
    bookmarks=true,colorlinks=true,linkcolor=blue,linktocpage,bookmarksopen=true,
    pdfauthor={\AuthorLong},
}

%% C-EDITINSTANCE-EXPAND: the kernel template colour keys want a literal colour
%% NAME and will NOT expand a macro — passing \ThemeAccent directly gives
%% "Unknown color '\ThemeAccent'", once per frame. \colorlet turns the macro
%% into a real named colour, so the indirection above still works.
\colorlet{themeaccent}{\ThemeAccent}
\colorlet{themeaccenttext}{\ThemeAccentText}

%% ---------------------------------------------------------------------------
%% Title page — ltx-talk's NATIVE \maketitle, restyled.
%%
%% ltx-talk ships a `titlepage' template with five `titlepage-element'
%% instances (title, subtitle, author, institute, date), each taking color,
%% font, before-skip, after-skip, tag-begin and tag-end. Restyling them gets a
%% designed title page while keeping \maketitle. The win is item 1 below; item 2
%% is what it costs nothing to keep, and is worth stating because it is the one
%% thing a hand-rolled frame also got right and would be easy to lose here:
%%
%%   1. ONE SOURCE OF TRUTH. Without \maketitle the long argument of \title is
%%      never rendered anywhere, so it silently duplicates the title frame's
%%      arguments and drifts from them. Only \title[short] stays live, consumed
%%      by the footer's element-order below. On one real course that had already
%%      produced a deck whose \title disagreed with its title page, and a deck
%%      with NO \title at all, whose footer was blank on every page with nothing
%%      to catch it.
%%   2. THE TITLE IS STILL THE DOCUMENT'S H1, and the class gives us a clean way
%%      to say so (see the tag-begin override below). A-HEADINGS depends on the
%%      deck having an H1 somewhere, and \section is otherwise the only one.
%%
%% It also removes the C-AND-TITLE hazard outright: the template does
%% \cs_set_protected:Npn \and {\quad} around the elements, so "A \and B" in
%% \author typesets correctly here with no local rebinding.
%%
%% ⚠ THE ATTRIBUTION STAYS WHERE BEAMER PUT IT. The beamer source writes the
%% "Material adapted from ..." block as free-standing centred text inside the
%% title frame, AFTER \maketitle, with \date{} empty:
%%
%%   \begin{frame}
%%     \maketitle
%%     \vspace{-1em}          %% re-tune: beamer's -4em is too much here
%%     \begin{center}
%%       Material adapted from: \\
%%       Russell and Norvig (AIMA Book): Chapters 7 and 9
%%     \end{center}
%%   \end{frame}
%%
%% That construct works VERBATIM under ltx-talk 0.6.2 -- measured, 0 errors,
%% Tagged: yes, nothing overlapping. It used to collide, which is what
%% C-TITLEPAGE was originally filed for, and that symptom is retired (#39). So
%% keep it: the deck's title block stays byte-identical to the beamer source and
%% switching a course back to beamer needs no edit here at all.
%%
%% Only the \vspace needs re-tuning. Beamer's -4em was compensating for beamer's
%% own title spacing; under ltx-talk it pulls the attribution up until it nearly
%% touches the institute line (measured: institute yMax 173.45, attribution
%% yMin 172.61). Set it by eye against the skips above.
%%
%% The `date' element below is restyled anyway, so a deck that prefers the
%% attribution inside the title block proper can put it in \date{...} instead.
%% That is a real alternative, not an error, but it costs something: it misuses a
%% semantic field on output whose whole point is accessibility, \@shortdate is
%% \let from \@date at class load so anything wanting a real short date gets the
%% attribution instead, a deck that wants BOTH a date and an attribution cannot
%% express it, and the title block stops matching the beamer source. See
%% references/compromises.md C-TITLEPAGE.
%%
%% ⚠ \date is NOT empty by default. \@date defaults to \today, so a deck that
%% sets no \date at all prints the DATE OF THE BUILD, changing on every rebuild.
%% The beamer decks already write \date{} -- keep it. convert_deck.py warns.
%%
%% ⚠ DO NOT leave a folded subtitle in \title. The beamer source writes
%%     \title[Short]{Main\\\large{Subtitle}}
%% and under the H1 override below that \\ makes the subtitle a SECOND H1 --
%% measured, /S /H1 = 2 against 1 when it is split. Split it, which \subtitle
%% does in beamer too, so the deck still compiles either way:
%%     \title[Short]{Main}
%%     \subtitle{Subtitle}
%% convert_deck.py reports an unsplit one.
%% ---------------------------------------------------------------------------

%% A-HEADINGS, on the title page. ltx-talk's stock `title' instance wraps the
%% title as \tag_struct_begin:n{tag=Title} #1 \par \tag_struct_end: -- a
%% MANUALLY opened struct around a paragraph that the kernel's AUTOMATIC
%% per-paragraph tagger also tags (as `text', RoleMapped to P). That nesting is
%% harmless while Title maps to plain P, and it is why you cannot simply remap
%% Title -> H1: that turns it into "<Hn> shall not contain <P>", which veraPDF's
%% PDF/UA-2 profile rejects. Same trap as C-CENTER-ARG, built into the class
%% rather than authored here.
%%
%% So do not retag the wrapper -- REPLACE it, and let the automatic tagger
%% produce the H1 directly, exactly as A-HEADINGS prescribes. tag-end is left
%% empty on purpose: the keys are scoped by the \group_begin:/\group_end: the
%% template already wraps each element in, so they revert on their own, and a
%% default \tag_struct_end: here would fire with no matching begin.
%%
%% Cost, stated: the PDF then has no /S /Title element. Nothing requires one --
%% the document title a checker reads is XMP dc:title, which comes from \title
%% and is unaffected (measured) -- and an H1 is worth more, because without it a
%% deck with no \section has no level-one heading at all.
\EditInstance{titlepage-element}{title}
  { color = themeaccent , font = \Huge\bfseries , before-skip = 0pt ,
    after-skip = 0.5em ,
    tag-begin = {\tagpdfsetup{para/tag=H1,para/flattened}} , tag-end = {} }
%% The stock subtitle is the blue `structure' colour; a typical beamer title
%% page has it black, so say so explicitly.
\EditInstance{titlepage-element}{subtitle}
  { color = black , font = \Large , before-skip = 0pt , after-skip = 2.5em }
\EditInstance{titlepage-element}{author}
  { font = \large , before-skip = 0pt , after-skip = 0.4em }
\EditInstance{titlepage-element}{institute}
  { font = \normalsize , before-skip = 0pt , after-skip = 2.5em }
\EditInstance{titlepage-element}{date}
  { font = \small , before-skip = 0pt , after-skip = 0pt }

%% ---------------------------------------------------------------------------
%% Compatibility shim.
%%
%%   \coursetitlepage{<main title>}{<subtitle>}{<attribution>}
%%
%% Earlier conversions hand-rolled the title frame with this macro, and it is
%% deployed across several courses. It now routes to the native title page, so
%% those decks need no edit and stop being able to drift: \title, the footer and
%% the printed title page all come from its first argument. A deck that never
%% called \title at all gets its footer title back as a side effect.
%%
%% NEW conversions should not use it. The beamer source already says what to
%% write -- \title/\subtitle/\date{} and \maketitle, with the attribution left as
%% centred text after it -- so there is nothing for a macro to package. The third
%% argument here goes to \date, which is the alternative form described above.
%%
%% The guard exists because a deck may never have called \title at all. On
%% ltx-talk 0.6.2 \@shorttitle is defined-and-empty from class load, so the
%% \ifx branch is the one that fires and \@ifundefined is belt and braces; on
%% 0.5.x \@shorttitle was genuinely UNDEFINED and \ifx alone took the wrong
%% branch, emitting an empty footer title in silence. \makeatletter is needed
%% because this file is \input, not a package.
%% ---------------------------------------------------------------------------
\makeatletter
\newcommand{\coursetitlepage}[3]{%
  \@ifundefined{@shorttitle}
    {\title{#1}}
    {\ifx\@shorttitle\@empty
       \title{#1}%
     \else
       \let\course@footername\@shorttitle
       \expandafter\title\expandafter[\course@footername]{#1}%
     \fi}%
  \subtitle{#2}%
  \date{#3}%
  \maketitle
}
\makeatother

%% ---------------------------------------------------------------------------
%% Visual style via ltx-talk templates
%% ---------------------------------------------------------------------------
\EditInstance{header}{std}{
  background-color = themeaccent,
  color            = themeaccenttext,
}
\EditInstance{footer}{std}{
  background-color = themeaccent,
  color            = themeaccenttext,
  element-order    = {author, title, framenumber},
  separator        = \hfill,
  left-hspace      = 3mm,
  right-hspace     = 3mm,
}

%% Tip: an EMPTY `background-color =' draws no bar at all — that is how you get
%% a minimalist, bar-less (Pittsburgh-like) header with a coloured frame title.

%% ---------------------------------------------------------------------------
%% C-TOC: Section dividers, without touching the decks.
%%
%% The old decks opened each section with an \AtBeginSection outline frame
%% (\tableofcontents[currentsection]). In ltx-talk 0.5.0 \tableofcontents
%% corrupts the PDF tag tree (mismatched para/text-unit hooks) and cannot be
%% produced as tagged output, so we cannot reproduce that outline as-is.
%%
%% Instead we REDEFINE \section itself to emit the section AND a clean,
%% tagging-safe divider frame. The decks keep their original \section{Title}
%% lines verbatim — no source edit, and no new macro to remember:
%%
%%   \section{Title}      section + divider frame   (the beamer habit)
%%   \section[Short]{Title}  ditto; Short goes to the TOC/navigation metadata
%%   \section*{Title}     section, NO divider frame (silent, escape hatch)
%%
%% \NewCommandCopy takes a snapshot of ltx-talk's own \section first, so the
%% redefinition delegates to it rather than recursing.
%% ---------------------------------------------------------------------------
\NewCommandCopy{\ltxtalkorigsection}{\section}
\RenewDocumentCommand{\section}{s o m}{%
  \IfBooleanTF{#1}{%
    %% starred: silent section, no divider
    \ltxtalkorigsection*{#3}%
  }{%
    \IfValueTF{#2}{\ltxtalkorigsection[#2]{#3}}{\ltxtalkorigsection{#3}}%
    \sectiondividerframe{#3}%
  }%
}

%% The divider frame itself — restyle this and every section page follows.
\newcommand{\sectiondividerframe}[1]{%
  \begin{frame}%
    \vfill
    \begin{center}
      {\color{\ThemeAccent}\Huge\bfseries #1}
    \end{center}
    \vfill
  \end{frame}%
}

%% Back-compat only: earlier versions of this skill told you to rewrite \section{X}
%% into \heading{X}. Decks converted then still say \heading. New decks should just
%% use \section.
\newcommand{\heading}[1]{\section{#1}}

%% ---------------------------------------------------------------------------
%% C-BACKGROUND: \usebackgroundtemplate, without touching the decks.
%%
%% ltx-talk has no \usebackgroundtemplate, so the decks' background frames stop
%% at "! Undefined control sequence". Stubbing it out is worse than the error:
%% these frames carry white text over a dark image, so the frame compiles, keeps
%% its page count, says Tagged: yes — and renders white on white.
%%
%% Define the command instead. The deck keeps its original lines verbatim,
%% including the { ... } group and the author's own \includegraphics call.
%%
%% Do NOT reach for an overlay tikz node (an earlier version of this skill said
%% to). A tikzpicture with [remember picture,overlay] is ordinary CONTENT that
%% draws outside its own box, so it paints in document order — over the header,
%% which is where ltx-talk draws the frame title. It also resolves current page
%% through the .aux, so any error that stops latexmk converging leaves the image
%% in the wrong place, silently. The kernel's shipout/background hook is the
%% layer ltx-talk uses for its own page colour (\__talk_pagecolor:n) and paints
%% beneath everything, which is what a background is.
%%
%% Three things the shim adds that the deck's own graphics call will not have:
%%
%%  * artifact — without it the background enters the structure tree as a
%%    /Figure whose /Alt is the FILENAME, once per SHIPPED page, so an n-overlay
%%    frame gets n bogus figures. (alt={} does NOT do this; it still emits the
%%    Figure with the filename as its /Alt.) Set through tag/graphic via a macro
%%    defined at top level: \ExplSyntaxOn does not survive being stored in hook
%%    code, so \keys_set:nn cannot be written inline inside \AddToHook.
%%  * height=\paperheight — beamer scaled by width and cropped the overflow.
%%    Under \put there is no crop, so width alone leaves a wider-than-page image
%%    short: measured 13% of the page WHITE for a 21:9 image on a 16:9 page,
%%    under white text. Full coverage wins; the cost is that a mismatched image
%%    is stretched rather than cropped.
%%  * \aftergroup — reproduces the scope the beamer original got from its
%%    { ... } group. Hook code is global, so without this the background would
%%    leak onto every later frame. It covers every overlay page of the frame;
%%    a single \AddToHookNext would cover only the first.
%%
%% Retire this whole block when ltx-talk gains a background-image interface of
%% its own: delete it, and the decks are already correct.
%% ---------------------------------------------------------------------------
\ExplSyntaxOn
\newcommand{\ltxtalkbgartifact}{\keys_set:nn{tag/graphic}{artifact}}
\ExplSyntaxOff
\newcommand{\ltxtalkbgoff}{\RemoveFromHook{shipout/background}[talkbg]}
\newcommand{\usebackgroundtemplate}[1]{%
  \AddToHook{shipout/background}[talkbg]{%
    \put(0cm,-\paperheight){%
      \begingroup\ltxtalkbgartifact\setkeys{Gin}{height=\paperheight}#1\endgroup}%
  }%
  \aftergroup\ltxtalkbgoff
}
